Theorem 1. Given a certain input x, let the causal graph in Fig. 1 (in the main paper) encode 2n causal patterns, i.e.,
Ω = 2N ={S :S⊆N} . If the causal effectwS of each causal patternS∈ Ω is measured by the Harsanyi dividend [29], i.e.
wS ≜∑
S′⊆S(−1)|S|−|S′|·v(xS′), then the causal graph faithfully encodes the inference logic of the DNN, as follows.
∀S⊆N , Y (xS) =v(xS) (10)
More crucially, the Harsanyi dividend is the unique metric that satisfies the faithfulness requirement.
• Proof: We only need to prove the following two statements. (1) Necessity: the causal graph based on Harsanyi dividends
wS satisfies the faithfulness requirement∀S⊆N ,Y (xS)=v(xS). (2) Sufficiency: if there exists another metric ˜wS that also
satisfies the faithfulness requirement, then, it is equivalent to the Harsanyi dividend, i.e.∀S⊆N , ˜wS =wS.
According to the SCM in Eq. (2) of the main paper, we haveY (xS) =∑
S′∈ΩwS′·CS′(xS) =∑
S′⊆SwS′. Therefore, the
faithfulness requirement can be equivalently re-written as∀S⊆N ,v (xS) =∑
S′⊆SwS′.
Proof for necessity. According to the definition of the Harsanyi dividend, we have∀S⊆N ,
14
